% ██████████████████████████████████████████████████████████████████
%  CHAPTER 2 — LITERATURE REVIEW
% ██████████████████████████████████████████████████████████████████
\chapter{Literature Review}
\label{ch:literature_review}

\section{Introduction}
\label{sec:lit_intro}

This chapter situates the thesis within four converging streams of literature: (i) multi-tenant SaaS architectures and the database-per-tenant isolation strategy adopted by modern ERP platforms; (ii) the application of tabular supervised learning to enterprise analytics; (iii) the evolution and theoretical foundations of post-hoc Explainable AI, with particular emphasis on SHAP and LIME; and (iv) the comparatively thin but growing body of work that interrogates the \emph{reliability} of explanations themselves. The chapter concludes by articulating the precise gap addressed by the present research.

\section{Multi-Tenant ERP Architectures}
\label{sec:lit_mt_erp}

The economics of enterprise software for SMEs are dominated by the multi-tenant SaaS delivery model \citep{bezemer2010multi, krebs2014multitenancy}. Three canonical isolation strategies---shared schema with row-level security, schema-per-tenant, and database-per-tenant---trade off operational complexity against the strength of data separation. This thesis adopts the database-per-tenant strategy implemented on Django~4.2 and PostgreSQL, motivated by its strong isolation guarantees and its compatibility with regulator-mandated data-residency requirements \citep{haddara2011erp}.

\section{Machine Learning in ERP Systems}
\label{sec:lit_ml_erp}

Embedded predictive services have become a baseline feature of modern ERP suites \citep{idc2024ai}. Tabular supervised learning algorithms---in particular gradient-boosted decision trees \citep{chen2016xgboost} and random forests---dominate this space due to their robustness on heterogeneous structured data, low feature-engineering burden, and competitive performance against deep architectures on tabular problems \citep{lundberg2020local}. Common deployment targets include late-payment prediction, lead scoring, and employee attrition---the three classifiers used as exemplars in this thesis.

\section{Explainable Artificial Intelligence}
\label{sec:lit_xai}

\subsection{The Black-Box Problem}
The opacity of high-capacity ML models has motivated a substantial literature on interpretability \citep{lipton2018mythos, rudin2019stop}. Two broad strategies have emerged: (i) intrinsically interpretable models, and (ii) post-hoc, model-agnostic explainers that operate after training. This thesis focuses exclusively on the latter, reflecting industry practice in ERP-embedded AI.

\subsection{SHAP}
SHAP unifies several earlier feature-attribution methods under the axiomatic framework of cooperative game theory, producing additive explanations whose components satisfy local accuracy, missingness, and consistency \citep{lundberg2017unified}. The \texttt{TreeExplainer} variant provides exact polynomial-time Shapley values for tree ensembles \citep{lundberg2020local}.

\subsection{LIME}
LIME approximates the local decision surface of a model around a target instance by fitting a sparse linear surrogate over perturbed samples \citep{ribeiro2016why}. It is model-agnostic by construction, but its outputs are sensitive to sampling and to the choice of locality bandwidth.

\section{Explanation Reliability and Failure Modes}
\label{sec:lit_xai_failure}

A smaller but growing body of work has begun to question whether XAI outputs are themselves trustworthy. Studies have documented instability under sampling perturbations, sensitivity to background-data selection, and adversarial manipulability of post-hoc explainers. Governance frameworks such as Model Cards \citep{mitchell2019model} and Datasheets for Datasets \citep{gebru2021datasheets} emphasize documentation but stop short of architectural controls. The forthcoming EU AI Act \citep{eu2024aiact} imposes transparency obligations but leaves implementation specifics to vendors. The present research extends this line of inquiry by quantifying explanation degradation specifically as a function of \emph{tenant data volume}, and by proposing an architectural---not procedural---mitigation.

\section{Research Gap}
\label{sec:lit_gap}

While prior work has separately addressed (i) multi-tenant ERP architecture, (ii) ML for enterprise analytics, and (iii) post-hoc XAI, no published study has, to the best of the author's knowledge, jointly characterised the degradation profile of SHAP and LIME across a controlled multi-tenant sparsity gradient \emph{and} proposed an integrated architectural safeguard. This thesis addresses that gap.
